\section{Introduction}
\label{sec:intro}

Most open language models are deployed with quantized weights. Storing weights in 4 or 8 bits instead of 16 cuts memory by $2$--$4\times$ and makes a 7B model fit on a laptop GPU. The cost of this compression is usually summarized by one number: a perplexity ratio on \wikitext, or accuracy on one multiple-choice benchmark~\cite{frantar2023gptq,lin2024awq,dettmers2023case}. People deploy these models for very different jobs, however, such as writing code, solving word problems, answering factual questions or drafting stories. If quantization damages these tasks unevenly, a single aggregate number can approve a bit width that silently breaks the one task a user cares about.

{\bf What do practitioners believe about quantization damage?} A common view is that INT8 is safe, INT4 is ``mostly fine'', and that tasks needing precise reasoning or memorized facts break first, possibly at a sharp cliff between INT8 and INT4. A second common view is that small models are more fragile than large ones. A third is that a distribution-level signal, such as the KL divergence between the quantized and full-precision next-token distributions, can flag trouble before task accuracy visibly drops. Each view is plausible, but none has been tested with a sweep fine enough to \emph{locate} a cliff across many task types.

{\bf What is missing?} Quantization papers such as \gptq~\cite{frantar2023gptq} and \awq~\cite{lin2024awq} report perplexity and a few zero-shot multiple-choice tasks, mostly at 3 and 4 bits. Later evaluation studies~\cite{li2024evaluating,jin2024comprehensive,marchisio2024quantization} show that degradation is uneven across tasks, and \citet{dutta2024accuracy} argue that accuracy hides ``flips'' and should be complemented by divergence metrics. To our knowledge, however, no study combines (i) a bit-width sweep fine enough to locate the cliff, (ii) generative tasks such as code, instruction following and creative writing alongside QA, and (iii) an explicit test of whether KL divergence \emph{leads} the accuracy cliff, with a numeric threshold.

{\bf Our approach.} We run a controlled sweep: seven instruction-tuned models from three families (\qwen 0.5B/1.5B/3B/7B, \llama 1B/8B, \phithree 3.8B), seven weight precisions (bf16, 8, 6, 5, 4, 3, 2 bits; round-to-nearest, group size 128), and six task types with 1,614 prompts in total. For every configuration we also measure the task-specific KL divergence from the bf16 model on the same task distribution. We pre-register four hypotheses. Arithmetic and factual QA should drop by more than 15\% between INT8 and INT4 (H1). Creative writing and reading comprehension should stay above 85\% at INT4 (H2). Smaller models should cliff at higher precision (H3). KL should cross about 0.3 nats at or before the cliff (H4). We complement the sweep with official \awq INT4 checkpoints and a per-channel quantization ablation.

{\bf What do we find?} Most of the hypotheses fail in their specific form, and the ways they fail are informative. The cliff for every task lies at or below 4$\to$3 bits: the largest drop is at 4$\to$3 bits in 28 of 30 cells for models up to 3.8B, and at 3$\to$2 bits in all 12 cells for the 7--8B models (\figref{fig:retained}). Eight and six bits are essentially lossless, with 3 significant drops in 72 paired model--task tests. At 4 bits, \emph{code} is the most sensitive task, not arithmetic: mean retained score is 0.84 for \humaneval versus 0.90 for \gsm and 0.97 for \ifeval. Size helps (Spearman $\rho=-0.53$, $p=0.0003$), but \qwen 3B is the least robust \qwen model at 3 bits. Task-specific KL detects broken configurations with AUC 0.97 and an optimal threshold of 0.21 nats (95\% CI 0.12--0.34), consistent with the hypothesized 0.3. It does not warn early, though: at the last healthy precision before a failure, KL exceeds 0.3 nats in only 6 of 42 cases. Finally, aggregate accuracy hides heavy answer churn, and more than 90\% of the answers lost at 4 bits are fluent and confident.

Our contributions are:
\begin{itemize}[leftmargin=*,itemsep=0pt,topsep=0pt]
    \item We map task-specific accuracy cliffs for 7 models $\times$ 7 precisions $\times$ 6 task types and show that the cliff lies between 4 and 3 bits for $\le$3.8B models and between 3 and 2 bits for 7--8B models, not between INT8 and INT4.
    \item We show that code generation is the most INT4-sensitive task and instruction following the least. We explain the ranking by whether correctness requires every token to be exact.
    \item We test KL divergence as an early-warning signal and find that it is an excellent \emph{concurrent} detector (AUC 0.97, threshold $\approx$0.2--0.3 nats) but not a \emph{leading} one, and that its useful threshold varies by about $100\times$ across tasks.
    \item We show that a fixed ``INT4'' label is not a fixed accuracy level. Removing the group-wise scales costs 30 points of retained score at 4 bits, while \awq recovers most of the round-to-nearest loss. Silent, fluent errors and answer churn are hidden by aggregate accuracy.
\end{itemize}

\Secref{sec:related} reviews related work, \secref{sec:method} describes the setup, \secref{sec:results} presents the results for each hypothesis, and \secref{sec:discussion} discusses implications and limitations.
